% step3-window-mixture.tex -- residual moment at BIM length, step 3 (window mixture).
% Fragment for v2/paper/main.tex (uses its macros \EE, \PP, \Var and its theorem environments).
% Algebra and exact identities: v2/scripts/residual/step3_window_mixture.py (sympy fallback;
% check tags [E1]-[E4], [A1]-[A5] below refer to that script).

\subsection*{Step 3: the window mixture for \(m\)-factors}

Summary. The plan was to put together the two host lemmas over the residual
degree window by the law of total variance. That plan fails at every \(m\ge1\), and
the failure is in the first moment. The conditional mean
\(\EE[Z\mid G,\deg R=g]\) is not flat across the window. The edge count
\(e(R)\) alone moves it by a factor \(e^{\Theta(\beta)}\) with
\(\beta^2=\Theta(m^2n/\delta)\to\infty\), and so the relative variance of
\(Z\) given \(G\) and the window \emph{diverges}
(Theorem~\ref{thm:s3-window-diverges}). This holds unconditionally and needs no
local limit theorem. What survives, given the host lemma of worker~2, is the
degree-conditioned statement, and for it no mixture and no LLT are needed.

\paragraph{Setting (S).}
\(n\) is even, \(G\) is a \(d\)-regular graph on \([n]\) with \(d\le n/2\),
\(C:=K_n\setminus G\) is \((n-1-d)\)-regular, and \(e_C:=e(C)=n(n-1-d)/2\).
Given \(G\), the residual \(R\) is the Bernoulli\((q)\) subgraph of \(C\)
(Proposition~\ref{prop:residual-degree-conditioned}), with mean degree
\(\delta:=(n-1-d)q\). The integer \(m\ge1\) has \(mn\) even, \(Z(H)\) is the number of
\(m\)-factors of \(H\), \(Z(C)>0\), and \(N:=mn/2\) is the number of edges of an
\(m\)-factor. Put
\[
k_0:=e_Cq=\frac{n\delta}{2},\qquad
\sigma^2:=e_Cq(1-q),\qquad
h(k):=\frac{(k)_N}{(e_C)_N},\qquad
\beta:=\frac{N\sigma}{k_0}=m\sqrt{\frac{n(1-q)}{2\delta}},
\]
\[
\eta_n:=\sqrt{\frac{6\log n}{\delta}},\qquad
\mathcal E_\eta:=\bigl\{\lvert\deg_R(v)-\delta\rvert\le\eta\delta\ \text{for all }v\bigr\},\qquad
W_\eta:=\bigl\{g:\lvert g_v-\delta\rvert\le\eta\delta\ \text{for all }v\bigr\}.
\]
\emph{BIM regime (B):} \(d/\log n\to\infty\), \(d=o\bigl(n^{3/8}(\log n)^{-1/3}\bigr)\),
\(\delta=\Theta\bigl(d^{2/3}(\log n)^{1/3}\bigr)\), \(m=\Theta\sqrt{d\log n}\).
In (B), \(m/\delta\to0\), \(\delta=o(n)\), \(\delta^3/n=\Theta(d^2\log n/n)\to0\), and
\[
\beta^2=\Theta\bigl(n\,d^{1/3}(\log n)^{2/3}\bigr),\qquad
\eta_n=\sqrt{6/C_3^3}\;\varepsilon_0\quad\bigl(\varepsilon_0=C_3(\log n/d)^{1/3}\bigr),\qquad
\frac{m}{\eta_n\delta}=\Theta\bigl(d^{1/6}(\log n)^{-1/6}\bigr)\to\infty
\]
(script [A5]). Note the last relation: an \(m\)-factor moves every degree by
\(m\), which is far outside the degree window.

\paragraph{Hypotheses taken from the parallel workers.}
Only (H2) is used. (H1) is its \(d=0\) case and is listed to fix the form.

\begin{hypothesis}[Form of lem:near-regular-host required]
\label{hyp:s3-near-regular}
In regime (B), with \(G=\emptyset\) (so \(C=K_n\)), the conclusion of
Hypothesis~\ref{hyp:s3-avoid-G} holds.
\end{hypothesis}

\begin{hypothesis}[Form of lem:avoid-G-host required]
\label{hyp:s3-avoid-G}
In regime (B) there is \(\varepsilon_n\to0\) such that the following holds for every \(d\)-regular \(G\) on \([n]\)
and every degree sequence \(g\in W_{\eta_n}\) with \(\sum_vg_v\) even and
\(\Omega(C,g)\neq\emptyset\). With \(H\sim\mathrm{Unif}(\Omega(C,g))\),
\(\EE_gZ:=\EE Z(H)>0\) and
\[
\Var_g Z\le\varepsilon_n\,(\EE_gZ)^2 .
\]
\end{hypothesis}

The window width needed is \(\eta_n=\sqrt{6\log n/\delta}\). In (B) this is a constant
multiple of \(\varepsilon_0\), so a host lemma proved on the residual-sharp window
\(\lvert g_v-\delta\rvert\le\varepsilon_0\delta/2\) suffices whenever
\(\sqrt{6/C_3^3}\le\tfrac12\). Otherwise the lemma is needed at width \(\eta_n\)
itself. Nothing about the dependence of \(\EE_gZ\) on \(g\) is assumed.

\begin{theorem}[Residual moment at BIM length: what holds and what fails]
\label{thm:residual-moment-bim-proof}
Assume setting (S) and regime (B).
\begin{enumerate}
\item[\textup{(i)}] \emph{(Window, proved.)} \(\PP(\deg R\in W_{\eta_n}\mid G)\ge1-2/n\).
\item[\textup{(ii)}] \emph{(Degree-conditioned moment, conditional on
Hypothesis~\ref{hyp:s3-avoid-G}.)} For every \(g\in W_{\eta_n}\) with
\(\PP(\deg R=g\mid G)>0\),
\[
\frac{\Var\bigl(Z(R)\bigm|G,\ \deg R=g\bigr)}{\EE\bigl[Z(R)\bigm|G,\ \deg R=g\bigr]^2}\le\varepsilon_n\to0 ,
\]
uniformly in \(G\) and \(g\). Consequently, for every \(t>0\),
\(\PP\bigl(\lvert Z(R)/\EE[Z(R)\mid G,\deg R]-1\rvert>t\bigm|G\bigr)\le\varepsilon_n/t^2+2/n\).
\item[\textup{(iii)}] \emph{(The window mixture fails, proved, no hypothesis.)} For every
\(\eta\in[\eta_n,\tfrac12]\),
\[
\frac{\Var\bigl(Z(R)\bigm|G,\ \mathcal E_\eta\bigr)}{\EE\bigl[Z(R)\bigm|G,\ \mathcal E_\eta\bigr]^2}\to\infty
\]
uniformly in \(G\), and the same holds with no window, conditioning on \(G\) alone.
\item[\textup{(iv)}] \emph{(The window means are not flat, conditional on~\cite[Theorem~4.6]{McK85}
as used in Theorem~\ref{thm:residual-matching-constrained}.)} The constant
sequences \(\delta\mathbf 1\) and \((\delta+2)\mathbf 1\) both lie in \(W_{\eta_n}\) when
\(\delta\in\mathbb Z\), and
\[
\frac{\EE_{(\delta+2)\mathbf 1}Z}{\EE_{\delta\mathbf 1}Z}
=\exp\Bigl(\frac{nm}{\delta}\bigl(1+O\bigl(\tfrac{m+1}{\delta}\bigr)\bigr)\Bigr)\to\infty .
\]
\end{enumerate}
In particular, Conjecture~\ref{conj:residual-moment-bim} is false as literally stated
(conditioning on \(G\) only) under the Bernoulli residual of
Proposition~\ref{prop:residual-degree-conditioned}. It is also false with the
residual-sharp window added. Its degree-conditioned form (ii) holds, conditional
on Hypothesis~\ref{hyp:s3-avoid-G}. All of this already applies at \(m=1\);
see Remark~\ref{rem:s3-m1}.
\end{theorem}

\begin{lemma}[Degree window, item (a)]
\label{lem:s3-window}
In setting (S), assume \(\delta\ge3\) and \(\eta_n\le\eta\le\frac12\).
\begin{enumerate}
\item[\textup{(a)}] \(\PP(\mathcal E_\eta^c\mid G)\le2/n\).
\item[\textup{(b)}] \emph{(Planted version.)} Let \(\lvert k-k_0\rvert\le n/2\) and \(k\ge N\). Let
\(F\) be a uniform \(m\)-factor of \(C\) and \(U\) a uniform \((k-N)\)-subset of
\(E(C)\setminus F\), and put \(\pi_k:=\PP(F\cup U\in\mathcal E_\eta)\). Then
\(\pi_k\ge1-3/n\).
\end{enumerate}
\end{lemma}

\begin{proof}
(a) \(\deg_R(v)\sim\mathrm{Bin}(n-1-d,q)\) has mean \(\delta\).
By~\cite[Corollary~2.3]{JLR00},
\(\PP(\lvert\deg_R(v)-\delta\rvert>\eta\delta)\le2e^{-\eta^2\delta/3}\le2e^{-\eta_n^2\delta/3}=2n^{-2}\).
Take a union over the \(n\) vertices.

(b) Fix \(F\). Then \(\deg_{F\cup U}(v)=m+X_v\), where \(X_v\) is hypergeometric:
\((k-N)\) draws from the \(n(n-1-d-m)/2\) edges of \(C\setminus F\), of which
\(n-1-d-m\) meet \(v\). So \(\mu:=\EE X_v=2(k-N)/n=2k/n-m\), and
\(\EE\deg_{F\cup U}(v)=2k/n\in[\delta-1,\delta+1]\) (script [E3] checks the mean
exactly on a small case). The window fails at \(v\) only if
\(\lvert X_v-\mu\rvert\ge t:=\eta\delta-1\). By~\cite[Theorem~2.10]{JLR00}
(Hoeffding~\cite{Hoeffding63}), \(X_v\) satisfies the binomial bounds
\(\PP(X_v\ge\mu+t)\le\exp(-t^2/(2(\mu+t/3)))\) and
\(\PP(X_v\le\mu-t)\le\exp(-t^2/(2\mu))\). Since \(\mu\le\delta+1\),
\(t\le\delta/2\) and \(\delta\ge3\), both bounds are at most \(\exp(-t^2/(3\delta))\). With
\(\eta\ge\eta_n\) we have \(t\ge\eta_n\delta-1\), and
\[
\frac{(\eta_n\delta-1)^2}{3\delta}=2\log n-\frac{2\eta_n}{3}+\frac1{3\delta}
\quad\text{(script [A3])},
\]
so each tail is at most \(e^{1/3}n^{-2}\). A union over \(v\) and both tails gives
\(1-\pi_k\le2e^{1/3}/n\le3/n\), uniformly in \(F\).
\end{proof}

\begin{lemma}[Exact edge-count identities]
\label{lem:s3-edge-count}
In setting (S), let \(\nu_k:=\PP(e(R)=k\mid G)\), so \(e(R)\sim\mathrm{Bin}(e_C,q)\).
Let \(\mathcal E\) be any set of subgraphs of \(C\) (for example \(\mathcal E_\eta\), or everything).
Let \(\pi_k\) be the planted probability of Lemma~\ref{lem:s3-window}(b) with
\(\mathcal E\) in place of \(\mathcal E_\eta\); set \(\pi_k\equiv1\) when \(\mathcal E\) is everything. Then
\[
\EE\bigl[Z(R)\mathbf 1_{\mathcal E}\bigm|G,\ e(R)=k\bigr]=Z(C)\,h(k)\,\pi_k ,
\qquad
\EE\bigl[Z(R)\bigm|G,\ e(R)=k\bigr]=Z(C)\,h(k).
\]
\end{lemma}

\begin{proof}
Given \(e(R)=k\), \(R\) is a uniform \(k\)-subset of \(E(C)\). Count pairs \((F,R)\)
with \(F\subseteq R\), \(\lvert R\rvert=k\), \(R\in\mathcal E\). Each of the \(Z(C)\)
factors \(F\) extends to \(\binom{e_C-N}{k-N}\) sets \(R=F\cup U\), and a fraction
\(\pi_k\) of those lie in \(\mathcal E\). Dividing by \(\binom{e_C}{k}\) and using
\(\binom{e_C-N}{k-N}/\binom{e_C}{k}=(k)_N/(e_C)_N\) gives the claim. Script [E1], [E2] verifies both identities by
exhaustive enumeration on \(C=K_6\setminus\)(perfect matching), \(m=1,2\), and several
windows. Script [E4] verifies the identity
\(\EE[Z\mathbf 1_{\mathcal E}\mid G]=Z(C)\sum_k\nu_kh(k)\pi_k\) as a polynomial in \(q\).
\end{proof}

\begin{theorem}[The window mixture diverges]
\label{thm:s3-window-diverges}
In setting (S), assume \(\delta\to\infty\), \(\delta=o(n)\), \(1\le m\le\delta/3\), and either
\(\mathcal E=\mathcal E_\eta\) with \(\eta_n\le\eta\le\tfrac12\), or \(\mathcal E\)
is everything. Let \(L=L_n\ge1\) satisfy
\[
L\le\frac{\beta}{10},\qquad 9L^2\le\log\sigma,\qquad L\le\frac16\sqrt{\frac n\delta}.
\]
Then, for all large \(n\) and uniformly in \(G\),
\[
\frac{\EE[Z^2\mid G,\mathcal E]}{\EE[Z\mid G,\mathcal E]^2}
\ \ge\ \Bigl(1-\frac2n\Bigr)\bigl(1-6e^{-L^2/2}\bigr)^2L^2 .
\]
All three upper bounds on \(L\) tend to infinity: \(\beta^2\ge n/(4\delta)\),
\(\sigma^2\ge n\delta/4\), and \(n/\delta\to\infty\). So the left side tends to \(\infty\).
\end{theorem}

\begin{proof}
Write \(f_k:=h(k)\pi_k\) and \(P:=\PP(\mathcal E\mid G)\). Then
\(\PP(e(R)=k\mid G,\mathcal E)=\nu_k\PP(\mathcal E\mid k)/P\). By
Lemma~\ref{lem:s3-edge-count}, \(\EE[Z\mid G,\mathcal E]=Z(C)\sum_k\nu_kf_k/P\) and
\(\EE[Z\mid G,e(R)=k,\mathcal E]=Z(C)f_k/\PP(\mathcal E\mid k)\) whenever
\(\PP(\mathcal E\mid k)>0\). If \(\PP(\mathcal E\mid k)=0\) then \(\pi_k=0\).
Conditional Jensen in \(k\) and \(\PP(\mathcal E\mid k)\le1\) give
\[
\EE[Z^2\mid G,\mathcal E]\ \ge\ \sum_k\frac{\nu_k\PP(\mathcal E\mid k)}{P}
\Bigl(\frac{Z(C)f_k}{\PP(\mathcal E\mid k)}\Bigr)^2
\ \ge\ \frac{Z(C)^2}{P}\sum_k\nu_kf_k^2
\quad\text{(script [E4] spot-checks this)}.
\]
Hence, with the tilted law \(w_k:=\nu_kf_k/\sum_j\nu_jf_j\) (script [A2]),
\[
\rho:=\frac{\EE[Z^2\mid G,\mathcal E]}{\EE[Z\mid G,\mathcal E]^2}
\ \ge\ P\,\frac{\sum_k\nu_kf_k^2}{(\sum_k\nu_kf_k)^2}
=P\sum_k\frac{w_k^2}{\nu_k}
\ \ge\ P\,\frac{w(A)^2}{\nu(A)}
\]
for every set \(A\), by Cauchy--Schwarz. Take \(A:=\{k\ge k_0+L\sigma\}\). Chebyshev gives
\(\nu(A)\le L^{-2}\), and \(P\ge1-2/n\) by Lemma~\ref{lem:s3-window}(a).
It remains to show \(w(A^c)\le6e^{-L^2/2}\).

Let \(J:=[k_0+2L\sigma,\,k_0+(2L+1)\sigma]\). Since \(3L\sigma\le n/2\) by the third
condition on \(L\) and \(\sigma^2\le n\delta/2\), every \(k\in J\) satisfies
\(\lvert k-k_0\rvert\le n/2\), and \(k\ge k_0\ge N\) because \(m\le\delta\). So \(\pi_k\ge1-3/n\) on \(J\)
(Lemma~\ref{lem:s3-window}(b)). For integers \(N\le k<k'\) one has
\(h(k)/h(k')=\prod_{i<N}(k-i)/(k'-i)\le(k/k')^N\), because
\((k-i)/(k'-i)-k/k'=-i(k'-k)/(k'(k'-i))\le0\) (script [A1]). Hence, for \(k<k_0+L\sigma\) and
\(k'\in J\), using \(\log(1+x)\ge x/(1+x)\) and \(2L\sigma\le k_0\),
\[
\frac{h(k)}{h(k')}\le\Bigl(\frac{k_0+L\sigma}{k_0+2L\sigma}\Bigr)^{N}
\le\exp\Bigl(-\frac{NL\sigma}{k_0+2L\sigma}\Bigr)\le e^{-\beta L/2}.
\]
The binomial Berry--Esseen bound
\(\sup_x\lvert\PP(e(R)\le x)-\Phi((x-k_0)/\sigma)\rvert\le0.56/\sigma\) and
\(9L^2\le\log\sigma\) give
\(\nu(J)\ge\phi(2L+1)-1.12/\sigma\ge e^{-9L^2/2}/(2\sqrt{2\pi})\) for large \(n\). Since \(h\) is increasing,
\[
w(A^c)\le\frac{\sum_{k<k_0+L\sigma}\nu_kh(k)}{\sum_{k\in J}\nu_kh(k)\pi_k}
\le\frac{e^{-\beta L/2}}{\nu(J)(1-3/n)}
\le\frac{2\sqrt{2\pi}}{1-3/n}\,e^{9L^2/2-\beta L/2}\le6e^{-L^2/2},
\]
because \(\beta L/2\ge5L^2\). Finally \(\beta^2=m^2n(1-q)/(2\delta)\) and
\(1-q\ge\tfrac12\), since \(q=\delta/(n-1-d)\le4\delta/n\to0\) when \(d\le n/2\). Also
\(\sigma^2=k_0(1-q)\ge n\delta/4\).
\end{proof}

\begin{proposition}[Window means are not flat, item (b)]
\label{prop:s3-means-not-flat}
Assume~\cite[Theorem~4.6]{McK85} in the form used in the proof of
Theorem~\ref{thm:residual-matching-constrained}. That form is: for degree sequences \(\mathbf h\)
with entries \(\le\delta+2\) avoiding a graph \(X\) of maximum degree \(\le d+m\),
\[
N(\mathbf h;X)=\frac{(S-1)!!}{\prod_vh_v!}
\exp\Bigl(-\lambda-\lambda^2-\sum_{uv\in X}\frac{h_uh_v}{S}+o(1)\Bigr),
\qquad S=\sum_vh_v,\quad \lambda=\sum_v\frac{h_v(h_v-1)}{2S},
\]
with \(o(1)\) uniform when \(\delta(d+\delta)^2/n\to0\), which holds in (B). In setting (S) with
\(\delta\in\mathbb Z\), \(n\delta\) even, and \(m\le\delta/3\),
\[
\frac{nm(\delta-2m-2)}{(\delta+2)(\delta-m)}-o(1)
\ \le\ \log\frac{\EE_{(\delta+2)\mathbf 1}Z}{\EE_{\delta\mathbf 1}Z}
\ \le\ nm\Bigl(\frac2{\delta-m}-\frac1\delta+\frac1{\delta^2}\Bigr)+o(1).
\]
\end{proposition}

\begin{proof}
For a constant sequence \(D\mathbf 1\) and any \(m\)-factor \(F\subseteq C\), deleting \(F\)
gives \(N(C,D\mathbf 1;F)=N((D-m)\mathbf 1;G\cup F)\). So
\(\EE_{D\mathbf 1}Z=Z(C)\exp(B(D)+\Delta(D)+o(1))\), independently of \(F\), where
\[
B(D)=n\log(D)_m+\log\frac{(n(D-m)-1)!!}{(nD-1)!!},\qquad
\Delta(D)=-f(D-m)+f(D)-\frac{(d+m)(D-m)}2+\frac{dD}2,
\]
with \(f(x)=\frac{x-1}2+\frac{(x-1)^2}4\). (The \(\lambda\) of a constant sequence \(x\) is
\((x-1)/2\), and \(\sum_{uv\in X}x^2/(nx)=rx/2\) for \(r\)-regular \(X\).) Then
\(\Delta(\delta+2)-\Delta(\delta)=0\) identically (script [A4]). For \(B\),
\(n\log\frac{(\delta+2)_m}{(\delta)_m}=n\sum_{j<m}\log\frac{\delta+2-j}{\delta-j}\) lies between
\(2nm/(\delta+2)\) and \(2nm/(\delta-m)\). The double-factorial part
\(-\sum_{i<N}\log\frac{n\delta+2n-1-2i}{n\delta-1-2i}\) lies between \(-nm/(\delta-m)\) and
\(-nm(1/\delta-1/\delta^2)\). Adding the bounds gives the display; script [A4] checks the
algebra and a numerical evaluation of \(B\) against both bounds. Both bounds are
\((nm/\delta)(1+O((m+1)/\delta))\). Finally \(2\le\eta_n\delta=\sqrt{6\delta\log n}\), so both
sequences lie in \(W_{\eta_n}\).
\end{proof}

\begin{remark}[Typical size of the non-flatness; heuristic]
\label{rem:s3-typical}
Expand \(\log\EE_gZ\) in \(y_v:=g_v-\delta\). The \(F\)-independent part \(B(g)\) has
linear coefficient
\(\kappa_1=\sum_{j<m}\frac1{\delta-j}-\sum_{i<N}\frac1{n\delta-1-2i}\approx\frac12\log\frac\delta{\delta-m}\approx\frac m{2\delta}\)
on \(S-n\delta=2(e(R)-k_0)\). It has quadratic coefficient
\(\kappa_2\approx m/(2\delta^2)\) on \(\sum_v(y_v-\bar y)^2\). Under the Bernoulli
law the two pieces have standard deviations \(\approx\beta\) and
\(\approx m\sqrt n/(\sqrt2\,\delta)\). In (B) these are
\(\Theta\bigl(\sqrt n\,d^{1/6}(\log n)^{1/3}\bigr)\) and
\(\Theta\bigl(\sqrt n\,d^{-1/6}(\log n)^{1/6}\bigr)\) (script [A5]). The heuristic value of
the relative variance in Theorem~\ref{thm:s3-window-diverges} is therefore
\(e^{\Theta(\beta^2)}=e^{\Theta(nd^{1/3}(\log n)^{2/3})}\), not just \(\omega(1)\). The second
piece shows that conditioning additionally on \(e(R)\), the \(G(n,M)\)-type residual,
would not rescue a mixture either: the degree spread alone moves the mean by
\(e^{\Theta(m\sqrt n/\delta)}\) whenever \(\delta=o(\sqrt n)\). This is the \(m\)-factor version of
the known \(G(n,M)\) threshold \(p\sqrt n\to\infty\) for perfect matchings.
\emph{Gap:} this remark is not proved. The Taylor remainder and the
\(F\)-dependent pairing sum \(\sum_{uv\in F}y_uy_v/S\) are uncontrolled at this
precision. Theorem~\ref{thm:s3-window-diverges} is the rigorous replacement.
\end{remark}

\begin{remark}[Consequences at \(m=1\)]
\label{rem:s3-m1}
Proposition~\ref{prop:s3-means-not-flat} and Theorem~\ref{thm:s3-window-diverges} hold for
\(m=1\). Proposition~\ref{prop:s3-means-not-flat} at \(m=1\) gives
\(\EE_{(\delta+2)\mathbf 1}Y/\EE_{\delta\mathbf 1}Y=e^{(1+O(1/\delta))n/\delta}\to\infty\). That
contradicts the uniform flatness asserted in Lemma~\ref{lem:residual-window-means}.
Its proof multiplies \(n/2\) per-edge factors \(1+O(\varepsilon_n)\), whose product is
not \(1+o(1)\). Theorem~\ref{thm:s3-window-diverges} at \(m=1\) gives
\(\Var(Y\mid G,\mathcal E_\eta)/\EE[Y\mid G,\mathcal E_\eta]^2\to\infty\). So
Corollary~\ref{cor:residual-matching-window}, the window form of
Theorem~\ref{thm:residual-moment}, and Corollary~\ref{thm:residual-moment-m1} fail whenever
the residual given \(G\) is Bernoulli on \(C\) with the window centred at its mean
degree. Theorem~\ref{thm:residual-matching-constrained} (exact degree \(\delta\)) is
unaffected, and so is the degree-conditioned form (ii) at \(m=1\) once that theorem
extends to \(g\in W_{\eta_n}\) (the extension already assumed in
Corollary~\ref{cor:residual-matching-window}). These statements in \texttt{main.tex}
were not edited here.
\end{remark}

\begin{remark}[Is a sparse multi-colour LLT needed?]
\label{rem:s3-llt}
No. A local limit theorem would control how the law of \(\deg R\) is spread
over the window. The obstruction is different: the conditional mean
\(\EE[Z\mid G,\deg R=g]\) varies by \(e^{\Theta(\beta)}\) across typical \(g\), and
no regularity of the degree law can flatten that. The only role of the window
is to put \(g\) where the host lemma applies (Lemma~\ref{lem:s3-window}(a),
Chernoff). The correct statement normalises by the random mean
\(\EE[Z\mid G,\deg R]\) (item (ii)), and it needs neither a mixture nor an LLT.
Proposition~\ref{prop:mllt-vanishing-margin} is therefore not on the critical path
of Conjecture~\ref{conj:residual-moment-bim}.
\end{remark}

\begin{proof}[Proof of Theorem~\ref{thm:residual-moment-bim-proof}]
(i) is Lemma~\ref{lem:s3-window}(a) with \(\eta=\eta_n\le\frac12\); in (B),
\(\eta_n\to0\) and \(\delta\to\infty\). (ii) By Theorem~\ref{thm:bern-degree-uniform},
\(\mathcal L(R\mid G,\deg R=g)=\mathrm{Unif}(\Omega(C,g))\), and
\(\Omega(C,g)\ne\emptyset\) because \(g\) is attained. So the ratio is
\(\Var_gZ/(\EE_gZ)^2\le\varepsilon_n\) by Hypothesis~\ref{hyp:s3-avoid-G}. The tail
bound is Chebyshev given \(\deg R\) on \(\{\deg R\in W_{\eta_n}\}\), plus (i). (iii) is
Theorem~\ref{thm:s3-window-diverges}: in (B), \(m\le\delta/3\), \(\delta=o(n)\) and
\(\beta,\sigma,n/\delta\to\infty\) (script [A5]). (iv) is
Proposition~\ref{prop:s3-means-not-flat}: in (B), \((m+1)/\delta\to0\), and
\(nm/\delta\ge n/\delta\to\infty\). The ``in particular'' follows: a relative
variance \(o(1)\) given \(G\) contradicts (iii) with no window. Given \(G\) and
the residual-sharp window \(\mathcal E_{\varepsilon_0/2}\), it contradicts (iii) with
\(\eta=\varepsilon_0/2\), which lies in \([\eta_n,\frac12]\) when \(\sqrt{6/C_3^3}\le\frac12\).
\end{proof}

\paragraph{Gaps and caveats.}
\begin{enumerate}
\item \emph{Hypothesis~\ref{hyp:s3-avoid-G}} (worker~2) is assumed, at window width
\(\eta_n=\sqrt{6\log n/\delta}\), uniformly in \(G\) and \(g\). Item (ii) is exactly as strong
as it is. Hypothesis~\ref{hyp:s3-near-regular} (worker~1) enters only through it.
\item \emph{Model mismatch in the existing text.}
Proposition~\ref{prop:residual-degree-conditioned} gives the residual density
\(p=(1+\varepsilon_0)d/n\) on \(C\). That has mean residual degree
\(\approx(1+\varepsilon_0)d\), not \(\delta=\varepsilon_0d\), so \(\mathcal E_\delta\) would then be a
large-deviation event. Setting (S) uses the density \(q\) whose mean degree is
\(\delta\). Item (ii) holds for every \(g\in W_{\eta_n}\) whatever the density, by
Theorem~\ref{thm:bern-degree-uniform}. Items (i) and (iii) use the centring. For an
off-centre window, (iii) is \emph{not proved}: the law of \(e(R)\) given the window is then
an exponential tilt piled at the window edge, and Lemma~\ref{lem:s3-window}(b) does not
apply as stated.
\item \emph{Downstream.} Proposition~\ref{prop:residual-to-peel-weight} asks for a
deterministic \(\nu_n\) with \(\EE[N_m\mid G]=(1+o(1))\nu_n\) and
\(\Var(N_m\mid G)=o(\nu_n^2)\). By (iii) this is impossible under the Bernoulli residual,
with or without the window, for every \(m\ge1\). Item (ii) supplies concentration only
around the random normaliser \(\EE[Z\mid G,\deg R]\). Whether the peel-weight and
total-variation argument of Conjecture~\ref{conj:config-split} can absorb a
\(g\)-dependent normaliser is \emph{open and not addressed here}. Item (iii) also does
not by itself show that \(N_m/\nu_n\not\to1\) in probability for every deterministic \(\nu_n\);
that would need a lower-tail bound on \(Z\) given \(e(R)\), and it is not proved.
\item Remark~\ref{rem:s3-typical} is heuristic. Item (iv) relies on~\cite[Theorem~4.6]{McK85}
exactly as the existing chain does. Items (i)--(iii) do not use McKay.
\end{enumerate}

\paragraph{Verification.}
\texttt{python v2/scripts/residual/step3\_window\_mixture.py} (sympy~1.14, the logged
fallback; no Wolfram app id). It runs:
[E1]--[E4], exhaustive enumeration of all \(2^{12}\) subgraphs of the octahedron
\(K_6\setminus\)(perfect matching) for \(m=1,2\) and four windows, checking
Lemma~\ref{lem:s3-edge-count}, the planted mean \(2k/n\), and the mixture identity and
inequality used in Theorem~\ref{thm:s3-window-diverges};
[A1] the two ratio inequalities;
[A2] the \(\chi^2\) identity;
[A3] the Chernoff exponent at \(\eta_n\);
[A4] \(\Delta(\delta+2)=\Delta(\delta)\), the bound algebra, and numerical values of
\(B(\delta+2)-B(\delta)\) inside the bounds (for example \(n=10^6,\delta=200,m=20\):
\(86909\le104260\le122722\), against \(nm/\delta=10^5\));
[A5] the regime-(B) exponents.
